\section{INTRODUCTION}

\subsection{Brief Overview}

Through this thesis project, I apply transformative game design theory and best practices to create a new competitive debating format. This research employs Research through Design (RtD) methodology to develop and playtest an alternative debate format that addresses what I believe to be systematic problems in widely-used formats, particularly World Schools and British Parliamentary debate.

\subsection{Motivation}

As a debater, judge, coach, and organizer with 15 years of experience, my insider perspective motivates this reformist work. Debate is a class of role-playing games with immense transformative potential, yet the design of major formats retains decades-old flaws, with game design theories entirely absent from attempts to improve them. 
\subsection{Problem Statement}

Although a very diverse range of debate formats are played all over the world, the iterative game-design improvement seen in other games is absent: most formats are deeply calcified, their competitive communities evolving new ``metas'', modes of play, and interpretations of rules without game-design supervision. As an example, I hypothesize that the two most popular debate formats, World Schools and British Parliamentary, whose communities in Europe are deeply intertwined, exhibit the following problematic tendencies:  Firstly, they seem to push debaters towards value convergence, an exclusionary and self-rightous holding of the same beliefs, potentially causing them to become less tolerant of people with different values and diminishing their ability to convince people from different backgrounds. Secondly, they appear to teach debaters to speak with extreme speed and information density, which could diminish their ability to meaningfully support their positions to people not trained as debate judges. Thirdly, in my opinion they sustain competitive pressure that erodes psychological safety and ossifies into status hierarchies among debaters.

These observations are not empirically validated, and this thesis does not intend to prove them, though a later section presents their echoes in literature.

\subsection{Research Question}

How can the application of transformative game design theory inform the creation of a competitive debate format that manages to elevate its transformative potential (improved critical thinking, improved discourse, self empowerment etc) by addressing the issues of value convergence, information density and competitive hierarchy?

This question decomposes in two: can a method of applying said theory to this problem be described and demonstrated? and does that method yield expected results?

\subsection{Scope and Limitations}

This thesis aims to demonstrate proof of concept: that alternative debate formats can be systematically designed to address problems, and that such designs produce observable, articulable differences in participant experience. The demonstration draws on participant observation and participant-reported experiences across five evaluation fields (engagement and appeal, productive pressure, growth direction, community dynamics, and epistemic humility), assessing whether a format grounded in transformative game design theory can produce clues of meaningful change without sacrificing the intellectual challenge or communal value of the activity.

The thesis does not aim to produce robust claims about long-term educational efficacy or entertainment value, as such claims would require longitudinal studies with larger samples than a master's thesis permits. Its main contribution to knowledge is to exemplify that debate formats can be designed better if looked at as roleplaying games.
